% ============================================================
%  Assignment 2: Architectural Choices, Tokenizers, Decoding, and Fine Tuning
% ============================================================
\documentclass[11pt, a4paper]{article}

\usepackage[utf8]{inputenc}
\usepackage[T1]{fontenc}
\usepackage[margin=2.5cm]{geometry}
\usepackage{booktabs}
\usepackage{array}
\usepackage{longtable}
\usepackage{amsmath}
\usepackage{hyperref}
\usepackage{xcolor}
\usepackage{enumitem}
\usepackage{float}
\usepackage{caption}

\hypersetup{
    colorlinks=true,
    linkcolor=blue!60!black,
    urlcolor=blue!60!black,
    citecolor=blue!60!black,
}

\setlength{\parskip}{4pt}
\setlength{\parindent}{0pt}

% Helper column types
\newcolumntype{L}[1]{>{\raggedright\arraybackslash}p{#1}}
\newcolumntype{R}[1]{>{\raggedleft\arraybackslash}p{#1}}
\newcolumntype{C}[1]{>{\centering\arraybackslash}p{#1}}

% SentencePiece word-start marker (U+2581 — not in standard TeX fonts)
\newcommand{\spmark}{\rule[0.18ex]{0.65em}{0.11ex}}

\newcommand{\src}[1]{\textsuperscript{\normalfont\small[#1]}}

\title{\textbf{Assignment 2}\\[6pt]
       \large Architectural Choices, Tokenizers, Decoding, and Fine Tuning}
\author{}
\date{May 2026}


\begin{document}
\maketitle
\tableofcontents
\newpage

% ============================================================
\section{Part 1 - Architectural Choices}
% ============================================================

% ----------------------------------------------------------
\subsection{Extraction Methodology}
% ----------------------------------------------------------

Architectural information for each model was gathered using \textbf{NotebookLM};
I ran a deep research experiment to collect sources for all models. 
After gathering a reliable list of sources, I asked \textbf{NotebookLM} to retrieve the relevant information for each model, while including references to the sources.

% ----------------------------------------------------------
\subsection{Per-Model Architectural Details}
% ----------------------------------------------------------

% ........................................................
\subsubsection{\texttt{meta-llama/Llama-3.1-8B-Instruct}}

\paragraph{1.~Layers and width.}
32 hidden layers\src{1}.
The hidden size is described as ``typically 4096--5120'' without specifying the
exact value for the 8B variant\src{1}; see Section~\ref{sec:uncertainty}.

\paragraph{2.~Attention heads.}
32 query heads per layer\src{1}.
The model uses Grouped-Query Attention (GQA)\src{1,2};
the exact number of KV heads is not reported in the available sources (see
Section~\ref{sec:uncertainty}).

\paragraph{3.~Component sizes.}
Dense decoder-only architecture; no MoE layers\src{1}.
Exact MLP intermediate size, attention head dimensions, and vocabulary and
embedding matrix sizes are all absent from the sources.

\paragraph{4.~Position encoding.}
RoPE\src{1}, extended via the \textit{llama3} multi-scale scheme\src{3}.
Context length: 128k tokens\src{2} (Source~1 states the tokenizer supports up to
8,192 tokens - disagreement noted in Section~\ref{sec:uncertainty}).
Scaling hyper-parameters: factor~=~8.0, low-freq factor~=~1.0,
high-freq factor~=~4.0, original max positions~=~8,192,
\texttt{rope\_type}~=~\texttt{llama3}\src{3}.

\paragraph{5.~Activations.}
SwiGLU, uniformly across all feed-forward sublayers\src{1}.

\paragraph{6.~Normalization.}
Pre-norm placement\src{1}.
Sources describe the norm type as ``RMSNorm or LayerNorm'' without definitively
specifying one\src{1}; RMSNorm is standard for this model family.

\paragraph{7.~Other properties.}
$\approx$8.03B total parameters\src{1}.
Instruction-tuned via SFT + DPO/RLHF\src{1}.
Supports 8 languages natively; knowledge cutoff December~2023\src{2}.
Training used 1.46M GPU-hours on H100-80\,GB hardware, producing
$\approx$420\,t\,CO\textsubscript{2}eq\src{2}.

\noindent\textbf{Sources:}
\begin{enumerate}[noitemsep,label=\small{[\arabic*]},leftmargin=*,topsep=2pt]
\item \small\url{https://www.emergentmind.com/topics/llama-3-1-8b-instruct-679bb9c6-400c-4f7c-b10a-c1048d659134}
\item \small\url{https://huggingface.co/meta-llama/Llama-3.1-8B-Instruct}
\item \small\url{https://huggingface.co/meta-llama/Llama-3.1-8B-Instruct/discussions/15}
\end{enumerate}

% ........................................................
\subsubsection{\texttt{mistralai/Mistral-7B-Instruct-v0.3}}

\paragraph{1.~Layers and width.}
32 hidden layers\src{2,3}; hidden size 4,096\src{2,3}.

\paragraph{2.~Attention heads.}
32 query heads, 8 KV heads (GQA, group size 4)\src{2,3}.
Faster inference via GQA is highlighted in\src{1,2}.

\paragraph{3.~Component sizes.}
Dense decoder; no MoE.
MLP intermediate size: 14,336 ($\approx3.5\times$ hidden)\src{2,3}.
Attention head dimension: 128\src{2}.
Vocabulary size: 32,768 (extended from 32,000 in v0.1/v0.2)\src{3,4}.
Separate (untied) input and output embedding matrices; exact sizes inferred from
hidden size and vocabulary.

\paragraph{4.~Position encoding.}
RoPE\src{2}; $\theta=1{,}000{,}000$; maximum context 32,768 tokens\src{3}.
Sliding Window Attention (window~=~4,096) is architecturally
supported\src{1,2}, but disabled (\texttt{sliding\_window: null}) in the v0.3
configuration\src{3}.

\paragraph{5.~Activations.}
SiLU, uniformly across all MLP layers\src{3}.

\paragraph{6.~Normalization.}
RMSNorm, $\varepsilon=10^{-5}$\src{3}.
Exact placement (pre- vs.\ post-norm) not stated in the sources; pre-norm is
the standard for this family.

\paragraph{7.~Other properties.}
$\approx$7.3B parameters\src{1}.
v0.3 release adds the v3 tokenizer and native function calling\src{4}.
Apache~2.0 license\src{1,3,4}.

\noindent\textbf{Sources:}
\begin{enumerate}[noitemsep,label=\small{[\arabic*]},leftmargin=*,topsep=2pt]
\item \small\url{https://mistral.ai/news/announcing-mistral-7b}
\item \small\url{https://www.emergentmind.com/papers/mistral-7b-instruct}
\item \small\url{https://huggingface.co/mistralai/Mistral-7B-Instruct-v0.3/blob/main/config.json}
\item \small\url{https://huggingface.co/mistralai/Mistral-7B-Instruct-v0.3}
\end{enumerate}

% ........................................................
\subsubsection{\texttt{Qwen/Qwen2.5-7B-Instruct}}

\paragraph{1.~Layers and width.}
28 hidden layers\src{1,2}; hidden size 3,584\src{2}.

\paragraph{2.~Attention heads.}
28 query heads, 4 KV heads (GQA, group size 7)\src{1,2}.

\paragraph{3.~Component sizes.}
Dense decoder; no MoE.
MLP intermediate size: 18,944 ($\approx5.3\times$ hidden) - notably wider than
peers\src{2}.
Attention head dimension not explicitly stated; inferred as 128 ($3584/28$).
Vocabulary size: 152,064\src{2}.
7.61B total parameters (6.53B non-embedding)\src{1}.
Untied input/output embeddings.

\paragraph{4.~Position encoding.}
RoPE with $\theta=1{,}000{,}000$\src{2}.
Base context: 32,768 tokens\src{2}; extended to 131,072 via
YaRN\src{1} (type=\texttt{yarn}, factor=4.0,
\texttt{original\_max\_position\_embeddings}=32,768\src{1}).
Generation is limited to 8,192 output tokens\src{1}.

\paragraph{5.~Activations.}
SiLU (\texttt{hidden\_act: silu}), uniformly\src{1,2}.

\paragraph{6.~Normalization.}
RMSNorm, $\varepsilon=10^{-6}$\src{1,2}.
Exact placement not stated in the sources.

\paragraph{7.~Other properties.}
Uses \textbf{QKV bias} in attention projections - uncommon among modern
models\src{1}.
Supports 29+ languages\src{1}.
Apache~2.0 license\src{1}.

\noindent\textbf{Sources:}
\begin{enumerate}[noitemsep,label=\small{[\arabic*]},leftmargin=*,topsep=2pt]
\item \small\url{https://huggingface.co/Qwen/Qwen2.5-7B-Instruct}
\item \small\url{https://huggingface.co/Qwen/Qwen2.5-7B-Instruct/blob/main/config.json}
\end{enumerate}

% ........................................................
\subsubsection{\texttt{allenai/OLMo-2-1124-7B-Instruct}}

\paragraph{1.~Layers and width.}
32 hidden layers; hidden size 4,096\src{1}.

\paragraph{2.~Attention heads.}
32 query heads and 32 KV heads - full Multi-Head Attention (MHA), no
KV-head reduction\src{1}.

\paragraph{3.~Component sizes.}
Dense decoder; no MoE.
MLP intermediate size: 11,008 ($\approx2.7\times$ hidden)\src{1}.
Head dimension inferred as 128 ($4096/32$).
Vocabulary size: 100,352; embeddings not tied (\texttt{tie\_word\_embeddings:
false})\src{1}.
No attention bias and no MLP bias\src{1}.

\paragraph{4.~Position encoding.}
RoPE, $\theta=500{,}000$\src{1}.
Maximum context: 4,096 tokens; \texttt{rope\_scaling} set to \texttt{null} -
no additional scaling applied\src{1}.

\paragraph{5.~Activations.}
SiLU, uniformly\src{1}.

\paragraph{6.~Normalization.}
RMSNorm, $\varepsilon=10^{-6}$\src{1}.
Exact placement not stated in the sources.

\paragraph{7.~Other properties.}
Post-trained via SFT + DPO + RLVR on the T-lu~3 dataset\src{2}.
Base model pre-trained on the Dolma dataset\src{2}.
All code, checkpoints, and training logs are publicly released\src{2}.

\noindent\textbf{Sources:}
\begin{enumerate}[noitemsep,label=\small{[\arabic*]},leftmargin=*,topsep=2pt]
\item \small\url{https://huggingface.co/allenai/OLMo-2-1124-7B-Instruct/blob/main/config.json}
\item \small\url{https://huggingface.co/eaddario/OLMo-2-1124-7B-Instruct-GGUF}
\end{enumerate}

% ........................................................
\subsubsection{\texttt{ibm-granite/granite-3.3-8b-instruct}}

\noindent\textit{Note: the instruct model shares the network geometry of its base
model \texttt{granite-3.3-8b-base}; parameters below are extracted from base
model documentation\src{1}.}

\paragraph{1.~Layers and width.}
40 hidden layers; hidden size 4,096\src{1}.

\paragraph{2.~Attention heads.}
32 query heads, 8 KV heads (GQA, group size 4)\src{1}.

\paragraph{3.~Component sizes.}
Dense decoder; no MoE\src{1}.
MLP intermediate size: 12,800 ($\approx3.1\times$ hidden)\src{1}.
Attention head dimension: 128\src{1}.
Shared input/output embeddings (tied)\src{1};
exact vocabulary size is not reported in the sources.

\paragraph{4.~Position encoding.}
RoPE; maximum context 128,000 tokens\src{1,2,3}.
Base $\theta$ and any scaling factors are not stated in the sources.

\paragraph{5.~Activations.}
SwiGLU, uniformly across all layers\src{1}.

\paragraph{6.~Normalization.}
RMSNorm\src{1}.
Exact placement and $\varepsilon$ not stated in the sources.

\paragraph{7.~Other properties.}
8.1B total parameters, all active at inference\src{1}.
Pretrained on 12T tokens\src{1}.
Natively supports \texttt{<think>}/\texttt{<response>} tags for structured
reasoning\src{2} and Fill-in-the-Middle (FIM) code completion\src{1}.
12 supported languages\src{2}.

\noindent\textbf{Sources:}
\begin{enumerate}[noitemsep,label=\small{[\arabic*]},leftmargin=*,topsep=2pt]
\item \small\url{https://huggingface.co/ibm-granite/granite-3.3-8b-base}
\item \small\url{https://build.nvidia.com/ibm/granite-3-3-8b-instruct}
\item \small\url{https://navigator.ufl.edu/}
\end{enumerate}

% ........................................................
\subsubsection{\texttt{deepseek-ai/DeepSeek-V3}}

\paragraph{1.~Layers and width.}
61 Transformer layers\src{2,3} (first 3 are dense FFN; the remaining 58 are
MoE)\src{3}.
Hidden dimension: 7,168\src{1,3}.

\paragraph{2.~Attention heads.}
128 query heads\src{3}.
The model uses Multi-Head Latent Attention (MLA) instead of standard MHA/GQA:
keys and values are compressed into a shared 512-dimensional latent vector
($d_c=512$), then decompressed per-head, dramatically reducing KV cache
memory\src{3,5}.
Query compression latent: $d_c'=1{,}536$\src{3}.
Attention head dimension: $d_h=128$\src{3}.

\paragraph{3.~Component sizes.}
Sparse MoE decoder.
Per MoE layer: 257 total experts - 1 shared expert (always active) plus 256
routed experts\src{2,3}.
Per token: 9 experts activated (1 shared + 8 routed)\src{2,3}.
Per-expert intermediate dimension: 2,048; combined active dimension:
$9\times2048=18{,}432$\src{1,3}.
Decoupled RoPE dimension per head: $d_h^R=64$\src{3}.
Shared input and output embedding matrices\src{3}.
Vocabulary size: 128,000\src{3} (Source~1 lists 129,280 - see
Section~\ref{sec:uncertainty}).

\paragraph{4.~Position encoding.}
Decoupled RoPE embedded within MLA: positional information is concatenated
to the KV representations rather than applied directly, resolving a mathematical
incompatibility with MLA compression\src{3,5}.
Maximum context: 131,072 tokens (128k)\src{3,4}.
Context extended from an initial 4k via YaRN: $s=40$, $\alpha=1$, $\beta=32$,
scaling factor $t=0.1\ln s+1$\src{3}.

\paragraph{5.~Activations.}
SwiGLU/SiLU, used consistently across all MoE and dense layers\src{1,2}.

\paragraph{6.~Normalization.}
RMSNorm, pre-norm placement\src{1,2,3}.
Additionally, extra RMSNorm layers are inserted \emph{after} the compressed MLA
latent vectors - a consequence of the MLA architecture\src{3}.

\paragraph{7.~Other properties.}
671B total parameters; 37B active per token (4.8\%)\src{2,3,4,6}.
Auxiliary-loss-free load balancing: a dynamic routing bias prevents collapse
without penalising the main loss\src{3,6}.
Multi-Token Prediction (MTP) training objective: predicts 2 tokens per step
(depth $D=1$)\src{3,5}.
First large-scale model pre-trained natively in FP8 mixed precision;
full training cost $\approx$\$5.58M on H800 GPUs\src{2,3,6}.

\noindent\textbf{Sources:}
\begin{enumerate}[noitemsep,label=\small{[\arabic*]},leftmargin=*,topsep=2pt]
\item \small\url{https://www.atlascloud.ai/blog/guides/analyzing-deepseek-v3-model-performance}
\item \small\url{https://fireworks.ai/blog/deepseek-model-architecture}
\item \small\url{https://arxiv.org/abs/2412.19437}
\item \small\url{https://github.com/deepseek-ai/DeepSeek-V3}
\item \small\url{https://mccormickml.com/2025/02/12/the-inner-workings-of-deep-seek-v3/}
\item \small\url{https://www.grandlinux.com/en/blogs/deepseek-moe-architecture.html}
\end{enumerate}

% ........................................................
\subsubsection{\texttt{HuggingFaceTB/SmolLM2-1.7B-Instruct}}

\paragraph{1.~Layers and width.}
24 hidden layers; hidden size 2,048\src{1}.

\paragraph{2.~Attention heads.}
32 query heads and 32 KV heads - full MHA, no KV reduction\src{1}.

\paragraph{3.~Component sizes.}
Dense decoder (\texttt{LlamaForCausalLM} architecture); no MoE\src{1}.
MLP intermediate size: 8,192 ($4\times$ hidden)\src{1}.
Head dimension inferred as 64 ($2048/32$).
Vocabulary size: 49,152\src{1,2}.
Tied word embeddings (\texttt{tie\_word\_embeddings: true})\src{1,2}.

\paragraph{4.~Position encoding.}
RoPE, $\theta=130{,}000$; \texttt{rope\_scaling} is \texttt{null} (no additional
scaling)\src{1}.
Maximum context: 8,192 tokens per \texttt{config.json}\src{1};
a secondary source lists 2,048 tokens\src{2} - see Section~\ref{sec:uncertainty}.

\paragraph{5.~Activations.}
SiLU (\texttt{hidden\_act: silu}), uniformly\src{1}.

\paragraph{6.~Normalization.}
RMSNorm, $\varepsilon=10^{-5}$\src{1}.
Exact placement not stated in the sources.

\paragraph{7.~Other properties.}
1.7B parameters\src{2,3}; designed for edge, mobile, and browser deployment.
Pretrained on 11T tokens (FineWeb-Edu, DCLM, The Stack, Cosmopedia v2, math
and coding datasets)\src{2,3}.
Post-trained via SFT + DPO on UltraFeedback\src{3}.
Training used 256 NVIDIA H100 GPUs in bfloat16\src{2,3}.

\noindent\textbf{Sources:}
\begin{enumerate}[noitemsep,label=\small{[\arabic*]},leftmargin=*,topsep=2pt]
\item \small\url{https://huggingface.co/HuggingFaceTB/SmolLM2-1.7B-Instruct/blob/main/config.json}
\item \small\url{https://collabnix.com/hugging-face-small-language-model-a-complete-guide/}
\item \small\url{https://mirai.us/models/SmolLM2-1.7B-Instruct}
\end{enumerate}

% ........................................................
\subsubsection{\texttt{microsoft/Phi-4-mini-instruct}}

\paragraph{1.~Layers and width.}
32 hidden layers; hidden size 3,072\src{2}.

\paragraph{2.~Attention heads.}
24 query heads, 8 KV heads (GQA, group size 3)\src{2,3,4}.

\paragraph{3.~Component sizes.}
Dense decoder; no MoE\src{1,3,4}.
MLP intermediate size: 8,192 ($\approx2.7\times$ hidden)\src{2}.
Head dimension inferred as 128 ($3072/24$).
Vocabulary size: 200,064\src{2,3,4}.
Tied word embeddings (\texttt{tie\_word\_embeddings: true})\src{1,2,3,4}.
No attention or MLP biases anywhere\src{2}.

\paragraph{4.~Position encoding.}
LongRoPE (\texttt{type: longrope})\src{2}.
Base $\theta=10{,}000$; original max positions: 4,096\src{2}.
Extended to 131,072 tokens\src{1,3,4} via a \texttt{rope\_scaling} dictionary
containing a \texttt{short\_factor} array (all 1.0) and a \texttt{long\_factor}
array with 48 dynamically scaled multipliers (ranging from 1.0 to 47.77)\src{2}.
Partial rotary factor: 0.75 - RoPE is applied to only 75\% of each head's
dimensions\src{2}.

\paragraph{5.~Activations.}
SiLU (\texttt{hidden\_act: silu}), uniformly\src{2}.

\paragraph{6.~Normalization.}
RMSNorm, $\varepsilon=10^{-5}$\src{2}.
Exact placement not stated in the sources.

\paragraph{7.~Other properties.}
3.8B parameters\src{1,3,4}.
Trained on 5T tokens (high-quality educational data, code, synthetic data,
filtered web text); knowledge cutoff June~2024\src{3,4}.
Post-trained via SFT + DPO + RLHF\src{1,3,4}.

\noindent\textbf{Sources:}
\begin{enumerate}[noitemsep,label=\small{[\arabic*]},leftmargin=*,topsep=2pt]
\item \small\url{https://catalog.models.ai.microsoft.com/}
\item \small\url{https://huggingface.co/microsoft/Phi-4-mini-instruct/blob/main/config.json}
\item \small\url{https://build.nvidia.com/microsoft/phi-4-mini-instruct}
\item \small\url{https://huggingface.co/microsoft/Phi-4-mini-instruct}
\end{enumerate}

% ........................................................
\subsubsection{\texttt{tiiuae/Falcon3-7B-Instruct}}

\paragraph{1.~Layers and width.}
28 hidden layers\src{1,2}; hidden size 3,072\src{1}.

\paragraph{2.~Attention heads.}
12 query heads, 4 KV heads (GQA, group size 3)\src{1,2}.
Attention head dimension: \textbf{256}\src{1,2} - double the 128 used by all
other models in this set.

\paragraph{3.~Component sizes.}
Dense decoder (\texttt{LlamaForCausalLM}); no MoE\src{1,2}.
MLP intermediate size: 23,040 ($\approx7.5\times$ hidden - largest ratio in this
set)\src{1}.
Vocabulary size: 131,072; untied embeddings
(\texttt{tie\_word\_embeddings: false})\src{1}.
No attention or MLP biases\src{1}.

\paragraph{4.~Position encoding.}
RoPE, $\theta=1{,}000{,}042$\src{1,2}; \texttt{rope\_scaling}~=~\texttt{null}\src{1}.
Maximum context: 32,768 tokens\src{1} (Source~2 summarises as ``32k'').

\paragraph{5.~Activations.}
SwiGLU (\texttt{hidden\_act: silu}), uniformly\src{1,2}.

\paragraph{6.~Normalization.}
RMSNorm, $\varepsilon=10^{-6}$\src{1,2}.
Exact placement not stated in the sources.

\paragraph{7.~Other properties.}
Despite the ``Falcon'' brand, the code architecture class is
\texttt{LlamaForCausalLM}\src{1}.
Pretrained on 14 teratokens; post-trained on 1.2M samples covering conversation,
code, STEM, safety, and function calling\src{2}.
Falcon-LLM-License (not Apache~2.0)\src{1,2}.

\noindent\textbf{Sources:}
\begin{enumerate}[noitemsep,label=\small{[\arabic*]},leftmargin=*,topsep=2pt]
\item \small\url{https://huggingface.co/tiiuae/Falcon3-7B-Instruct/blob/main/config.json}
\item \small\url{https://build.nvidia.com/tiiuae/falcon3-7b-instruct}
\end{enumerate}

% ........................................................
\subsubsection{\texttt{dicta-il/dictalm2.0-instruct}}

\paragraph{1.~Layers and width.}
32 hidden layers; hidden size 4,096\src{2}.

\paragraph{2.~Attention heads.}
32 query heads, 8 KV heads (GQA, group size 4)\src{2}.

\paragraph{3.~Component sizes.}
Dense decoder based on Mistral-7B-v0.1; no MoE\src{2,3}.
MLP intermediate size: 14,336 ($\approx3.5\times$ hidden)\src{2}.
Attention head dimension (RoPE dimension count): 128\src{2}.
Vocabulary size: 33,152\src{2}.
Separate (untied) input embedding and output matrices, each of size
$4096\times33{,}152$\src{2}.

\paragraph{4.~Position encoding.}
RoPE; base frequency ($\theta$): 10,000; applied over 128 dimensions\src{2}.
Maximum context: 32,768 tokens\src{1,2}.

\paragraph{5.~Activations.}
Activation function not explicitly named in the sources.
Inferred as SiLU/SwiGLU from the presence of \texttt{ffn\_gate},
\texttt{ffn\_up}, and \texttt{ffn\_down} weight matrices\src{2}, consistent with
the Mistral-7B-v0.1 base architecture\src{3}.

\paragraph{6.~Normalization.}
RMSNorm, $\varepsilon=10^{-5}$\src{2}.
Pre-norm placement is confirmed by the presence of \texttt{attn\_norm},
\texttt{ffn\_norm}, and \texttt{output\_norm} weight tensors in the model
blobs\src{2}.

\paragraph{7.~Other properties.}
The only Hebrew-specialised model in this set.
Built on Mistral-7B-v0.1, then continually pre-trained on 190B tokens
(50\% Hebrew, 50\% English)\src{3}.
Hebrew tokenisation rate improved from $\approx$5.78 to \textbf{2.76
tokens/word} through the injection of $\approx$1,000 Hebrew-specific tokens into
the vocabulary\src{3}.

\noindent\textbf{Sources:}
\begin{enumerate}[noitemsep,label=\small{[\arabic*]},leftmargin=*,topsep=2pt]
\item \small\url{https://ollama.com/aminadaven/dictalm2.0-instruct}
\item \small\url{https://ollama.com/aminadaven/dictalm2.0-instruct:q4_0}
\item \small\url{https://huggingface.co/dicta-il/dictalm2.0}
\end{enumerate}

% ----------------------------------------------------------
\subsection{Uncertainties, Missing Fields, and Source Disagreements}
\label{sec:uncertainty}
% ----------------------------------------------------------

\begin{description}[leftmargin=2.2cm,labelwidth=2cm,style=sameline]

\item[\textbf{Llama}]
  Sources\src{1} describe the hidden size as ``typically 4096--5120'' without
  pinning the 8B model to a specific value; 4096 is used as the well-known
  value for this variant.
  KV head count is absent from all three sources; 8 is the accepted value.
  Norm type is listed as ``RMSNorm \emph{or} LayerNorm''\src{1} without
  resolution.
  Source~1 states the tokenizer supports 8,192 tokens, while Source~2 lists the
  context length as 128k; the latter is treated as authoritative for the
  \texttt{max\_position\_embeddings} configuration.

\item[\textbf{DeepSeek-V3}]
  Vocabulary size: the arXiv technical report\src{3} states 128,000, while an
  Atlas Cloud analysis\src{1} lists 129,280.
  Source~3 is treated as authoritative.
  The MLP ``size'' of 2,048 refers to each individual expert's intermediate
  dimension; the effective combined dimension across all 9 active experts is
  18,432, which is the figure Source~1 reports as the MLP dimension.

\item[\textbf{Granite}]
  Vocabulary size is not reported in any of the three sources consulted.
  RoPE $\theta$ value and any scaling parameters are also absent; these fields
  are left as NA.

\item[\textbf{SmolLM2}]
  Source~2 states the context length as 2,048 tokens, while
  \texttt{config.json}\src{1} sets \texttt{max\_position\_embeddings} to 8,192.
  The configuration file is treated as authoritative; the discrepancy likely
  reflects documentation lag or a deployment-level truncation.

\end{description}

% ----------------------------------------------------------
\subsection{Architectural Comparison}
% ----------------------------------------------------------

\begin{longtable}{@{}L{3.0cm}rrrrrrl@{}}
\caption{Architectural parameters - all 10 models}\label{tab:comparison}\\
\toprule
\textbf{Model} & \textbf{Lay.} & \textbf{Hidden} & \textbf{Q-h} & \textbf{KV-h} & \textbf{MLP} & \textbf{Ctx\,(k)} & \textbf{Vocab} \\
\midrule
\endfirsthead
\multicolumn{8}{l}{\small\textit{(continued)}}\\
\toprule
\textbf{Model} & \textbf{Lay.} & \textbf{Hidden} & \textbf{Q-h} & \textbf{KV-h} & \textbf{MLP} & \textbf{Ctx\,(k)} & \textbf{Vocab} \\
\midrule
\endhead
\bottomrule
\endfoot
Llama-3.1-8B    & 32 & 4096 & 32 & 8  & 14336 & 128 & 128k \\
Mistral-7B v0.3 & 32 & 4096 & 32 & 8  & 14336 & 32  & 33k  \\
Qwen2.5-7B      & 28 & 3584 & 28 & 4  & 18944 & 128 & 152k \\
OLMo-2-7B       & 32 & 4096 & 32 & 32 & 11008 & 4   & 100k \\
Granite-3.3-8B  & 40 & 4096 & 32 & 8  & 12800 & 128 & NA   \\
DeepSeek-V3     & 61 & 7168 & 128 & MLA & 2048 (per~exp.) & 128 & 128k \\
SmolLM2-1.7B    & 24 & 2048 & 32 & 32 & 8192  & 8   & 49k  \\
Phi-4-mini      & 32 & 3072 & 24 & 8  & 8192  & 128 & 200k \\
Falcon3-7B      & 28 & 3072 & 12 & 4  & 23040 & 32  & 131k \\
DictaLM2.0      & 32 & 4096 & 32 & 8  & 14336 & 32  & 33k  \\
\end{longtable}

\begin{table}[H]
\centering\small
\caption{Activation, norm, and position-encoding scheme per model}
\label{tab:comparison2}
\begin{tabular}{@{}L{3.0cm}L{1.8cm}L{1.8cm}L{7cm}@{}}
\toprule
\textbf{Model} & \textbf{Act.} & \textbf{Norm} & \textbf{Position encoding} \\
\midrule
Llama-3.1-8B    & SwiGLU & RMSNorm & RoPE, llama3 multi-scale scaling \\
Mistral-7B v0.3 & SiLU   & RMSNorm & RoPE, $\theta$=1M; SWA disabled \\
Qwen2.5-7B      & SiLU   & RMSNorm & RoPE + YaRN (factor 4.0) \\
OLMo-2-7B       & SiLU   & RMSNorm & RoPE, $\theta$=500k, no scaling \\
Granite-3.3-8B  & SwiGLU & RMSNorm & RoPE, 128k ctx (params unknown) \\
DeepSeek-V3     & SwiGLU & RMSNorm & Decoupled RoPE + YaRN in MLA \\
SmolLM2-1.7B    & SiLU   & RMSNorm & RoPE, $\theta$=130k, no scaling \\
Phi-4-mini      & SiLU   & RMSNorm & LongRoPE, partial rotation (0.75) \\
Falcon3-7B      & SwiGLU & RMSNorm & RoPE, $\theta$=1M, head\_dim=256 \\
DictaLM2.0      & SiLU   & RMSNorm & RoPE, $\theta$=10k (lowest in set) \\
\bottomrule
\end{tabular}
\end{table}

\paragraph{Attention mechanism.}
Seven of ten models use GQA (predominantly 8 KV heads for a 4:1 ratio against 32
Q heads).
Two models - OLMo-2 and SmolLM2 - retain full MHA (32 KV heads = 32 Q heads),
which is viable at their scale where KV cache pressure is lower.
DeepSeek-V3 uses MLA, eliminating traditional KV heads by compressing into a
512-dim latent.

\paragraph{MLP width.}
The intermediate-to-hidden ratio spans from $2.7\times$ (OLMo-2, Phi-4-mini)
to $7.5\times$ (Falcon3).
The majority cluster at $3.5\times$ (Llama, Mistral, DictaLM).
Qwen ($5.3\times$) and Falcon ($7.5\times$) are wide outliers;
Falcon compensates for a narrow hidden size and few Q heads by widening the FFN.

\paragraph{Context length strategies.}
All models use RoPE but differ in how they extend it: a high base $\theta$ alone
(Mistral: 1M, Falcon: $\approx$1M), YaRN on top of a high $\theta$ (Qwen, DeepSeek),
llama3 multi-scale frequency stretching (Llama), and LongRoPE interpolation (Phi).
DictaLM's $\theta$=10,000 is the pre-scaling-era default, making it the weakest
at long-context generalisation.

\paragraph{Vocabulary size.}
Ranges from 33k (Mistral, DictaLM) to 200k (Phi-4-mini).
Larger vocabularies reduce tokenisation overhead for code and non-English text
at the cost of larger embedding matrices.

% ----------------------------------------------------------
\subsection{Reflection and Analysis}
% ----------------------------------------------------------

\paragraph{Consensus choices.}
RMSNorm with pre-norm, RoPE positional encoding, and a gated linear unit
activation (SiLU or SwiGLU) are universal
across all ten models.
GQA is near-universal for the 7--8B class (7 of 8 applicable models).

\paragraph{Outliers and rules of thumb.}
\begin{itemize}[noitemsep]
  \item \textbf{Falcon3:} FFN ratio $7.5\times$; head dimension 256 (vs.\ 128
        everywhere else); fewest Q heads (12). A ``wide-head'' design that
        compensates narrowness in other dimensions with an unusually large FFN.
  \item \textbf{OLMo-2:} full MHA, $2.7\times$ FFN, 4k context. Prioritises
        transparency and reproducibility over deployment efficiency.
  \item \textbf{Qwen2.5:} $5.3\times$ FFN, QKV bias, 152k vocab. Heavy
        investment in multilingual and reasoning performance.
  \item \textbf{DeepSeek-V3:} qualitatively different class - MoE with 671B
        total / 37B active, MLA attention, FP8 training.
  \item \textbf{DictaLM2.0:} $\theta$=10,000 is the pre-high-$\theta$ Mistral
        default, leaving it the least robust for long-sequence inputs.
  \item The most common Q/KV ratio among GQA models is 4:1 (32Q/8KV at
        hidden=4096). Qwen uses 7:1; Falcon and Phi use 3:1.
  \item Deeper stack (40 layers) for Granite; MoE depth (61 layers) for
        DeepSeek, where active parameter count per forward pass remains
        tractable.
\end{itemize}

\paragraph{Recommended design for a new 7--8B model.}
Based on the consensus above: 32 layers, hidden size 4,096, GQA with 32Q/8KV
heads and head dimension 128.
SwiGLU MLP with intermediate size $\approx14\text{k}$--$16\text{k}$
($3.5$--$4\times$ hidden).
RoPE with $\theta\approx500\text{k}$--$1\text{M}$; YaRN can be applied at
inference time if contexts $>32\text{k}$ are needed, avoiding baking in complex
scaling during pretraining.
RMSNorm pre-norm, $\varepsilon=10^{-5}$.
Vocab $\approx100\text{k}$--$130\text{k}$ for broad multilingual coverage;
tied embeddings only if the model is small or memory-constrained.
Avoid MoE and MLA unless the deployment scenario explicitly requires them.

% ============================================================
\newpage
\section{Part 2 - Tokenizers}
% ============================================================

% ----------------------------------------------------------
\subsection{Tokenizer Choices per Model}
% ----------------------------------------------------------

Table~\ref{tab:tokenizers} summarises the three tokenizer properties examined
for each model: vocabulary size, word-boundary strategy, and non-word tokens.

\begin{table}[H]
\centering
\caption{Tokenizer summary for all 10 models}
\label{tab:tokenizers}
\small
\begin{tabular}{L{3.4cm} L{3.6cm} R{1.3cm} L{2.4cm} C{1.1cm} R{1.1cm}}
\toprule
\textbf{Model} & \textbf{Class} & \textbf{Vocab} & \textbf{Boundary} & \textbf{Byte} & \textbf{Special} \\
\midrule
Llama-3.1-8B    & \texttt{PreTrainedTokenizerFast} & 128{,}256 & \texttt{Ġ} (BPE)          & No  &  2 \\
Mistral-7B      & \texttt{LlamaTokenizerFast}      &  32{,}768 & \texttt{\spmark} (SentencePiece)& Yes &  3 \\
Qwen2.5-7B      & \texttt{Qwen2TokenizerFast}      & 151{,}665 & \texttt{Ġ} (BPE)          & No  & 14 \\
OLMo-2-7B       & \texttt{GPT2TokenizerFast}       & 100{,}278 & \texttt{Ġ} (BPE)          & No  &  2 \\
Granite-3.3-8B  & \texttt{GPT2TokenizerFast}       &  49{,}159 & \texttt{Ġ} (BPE)          & No  &  8 \\
DeepSeek-V3     & \texttt{LlamaTokenizerFast}      & 128{,}815 & \texttt{Ġ} (BPE)          & No  &  2 \\
SmolLM2-1.7B    & \texttt{GPT2TokenizerFast}       &  49{,}152 & \texttt{Ġ} (BPE)          & No  &  3 \\
Phi-4-mini      & \texttt{GPT2TokenizerFast}       & 200{,}029 & \texttt{Ġ} (BPE)          & No  &  1 \\
Falcon3-7B      & \texttt{PreTrainedTokenizerFast} & 131{,}072 & \texttt{Ġ} (BPE)          & No  & 25 \\
DictaLM-2.0     & \texttt{LlamaTokenizerFast}      &  33{,}152 & \texttt{\spmark} (SentencePiece)& Yes &  3 \\
\bottomrule
\end{tabular}
\end{table}

\paragraph{Vocabulary size.}
Vocabulary sizes span nearly an order of magnitude, from 32,768 tokens
(Mistral-7B, DictaLM-2.0) to 200,029 (Phi-4-mini), reflecting different
trade-offs between coverage, embedding-table size, and average sequence length.

\paragraph{Word-boundary strategy.}
All ten tokenizers use a \emph{prefix-marker} convention to encode which tokens
are the beginning of a new word, so that the original text can be reconstructed
unambiguously from the token sequence alone.  Eight models use
\textbf{GPT-2-style BPE}: a leading \texttt{Ġ} (U+0120) signals that the token
opens a new word; tokens without this prefix are continuations that attach
directly to the preceding token.  The two exceptions---\textbf{Mistral-7B} and
\textbf{DictaLM-2.0}---use \textbf{SentencePiece}, which applies the same
logic with \texttt{\spmark} (U+2581) instead.  Reconstruction is thus
straightforward in both cases: tokens are concatenated left-to-right; each
occurrence of the boundary marker (stripped from the token string and replaced
by a space) starts a new word.

\paragraph{Special tokens.}
The ``Special'' column counts tokens that encode structural or control
information rather than text content: sequence delimiters (\texttt{<s>},
\texttt{</s>}, \texttt{<|endoftext|>}), unknown-token placeholders
(\texttt{<unk>}), and chat-format markers.  Most models have 1--3 such tokens.
Three models stand out:
\begin{itemize}[noitemsep]
  \item \textbf{Qwen2.5} (14): adds vision/tool-use delimiters
        (\texttt{<|im\_start|>}, \texttt{<|im\_end|>},
        \texttt{<|object\_ref\_start|>}, \texttt{<|box\_start|>},~\ldots).
  \item \textbf{Granite-3.3} (8): adds role-marking tokens
        (\texttt{<|start\_of\_role|>}, \texttt{<|end\_of\_role|>},
        \texttt{<|tool\_call|>}, \texttt{<|start\_of\_cite|>},~\ldots).
  \item \textbf{Falcon3} (25): inherits legacy corpus-metadata markers from
        pre-training (\texttt{>>TITLE<<}, \texttt{>>ABSTRACT<<},~\ldots).
\end{itemize}

\paragraph{Byte-fallback tokens.}
Only the two SentencePiece tokenizers (Mistral, DictaLM) carry an additional
256 \texttt{<0xNN>} byte tokens, one per possible byte value.  These are also
non-word tokens: they do not represent text directly but serve as a universal
fallback, encoding any Unicode codepoint byte-by-byte whenever it lacks a
dedicated vocabulary entry.  This guarantees that the tokenizer is losslessly
invertible on arbitrary input.

% ----------------------------------------------------------
\subsection{Average Tokens per Word}
% ----------------------------------------------------------

\paragraph{Methodology.}
Two reference corpora were used: a $\sim$1{,}000-word English text and a
$\sim$1{,}000-word Hebrew text (both stored in \texttt{docs/}).
For each model, the corpus was encoded with
\texttt{tok.encode(text, add\_special\_tokens=False)}; the \emph{word count}
was defined as the number of whitespace-delimited tokens
(\texttt{len(text.split())}).  Average tokens per word is therefore
$\frac{\text{encoded token count}}{\text{whitespace word count}}$.

This method assumes that whitespace splitting is a reasonable word-count proxy
for both English and Hebrew.  Hebrew words are whitespace-delimited (unlike,
e.g., Chinese), so the same formula applies directly.  The ratio does not
account for punctuation attached to words, but this bias is uniform across
models and therefore does not affect \emph{comparisons}.

\begin{table}[H]
\centering
\caption{Average tokens per word on English and Hebrew corpora}
\label{tab:avgtokens}
\small
\begin{tabular}{L{3.4cm} R{2.2cm} R{2.2cm} R{2.0cm}}
\toprule
\textbf{Model} & \textbf{EN avg} & \textbf{HE avg} & \textbf{HE/EN ratio} \\
\midrule
Llama-3.1-8B    & 1.29 & 5.77 & 4.47 \\
Mistral-7B      & 1.49 & 5.97 & 4.01 \\
Qwen2.5-7B      & 1.33 & 2.58 & 1.94 \\
OLMo-2-7B       & 1.30 & 5.78 & 4.45 \\
Granite-3.3-8B  & 1.76 & 5.28 & 3.00 \\
DeepSeek-V3     & 1.28 & 2.86 & 2.23 \\
SmolLM2-1.7B    & 1.33 & 6.39 & 4.80 \\
Phi-4-mini      & 1.29 & 2.42 & 1.88 \\
Falcon3-7B      & 1.30 & 5.31 & 4.08 \\
DictaLM-2.0     & 1.49 & 2.84 & 1.91 \\
\bottomrule
\end{tabular}
\end{table}

% ----------------------------------------------------------
\subsection{Tokenization Differences Example}
% ----------------------------------------------------------

\paragraph{Text used.}
The following 59-word English passage was used
(file \texttt{docs/english\_text\_part2.txt}):

\begin{quote}
\small
The World Chess Championship 2026 is an upcoming chess match to determine the
World Chess Champion.  It will be played between the defending champion Gukesh
Dommaraju and the challenger, Javokhir Sindarov, the winner of the Candidates
Tournament 2026.  The Championship is provisionally scheduled between 23
November and 17 December 2026, while the host city is yet to be decided.
\end{quote}

The span \textbf{\texttt{"Gukesh Dommaraju"}} (rare South-Indian proper name)
was chosen because it produces the greatest variation across tokenizers.
Table~\ref{tab:gukesh} shows the tokenization by all ten models; four distinct
segmentations are produced.

\begin{table}[H]
\centering
\caption{Tokenization of \texttt{"Gukesh Dommaraju"} across all 10 models}
\label{tab:gukesh}
\small
\begin{tabular}{L{3.4cm} L{9cm}}
\toprule
\textbf{Model(s)} & \textbf{Tokens} \\
\midrule
Llama, Qwen, OLMo, DeepSeek, Phi, Falcon
  & \texttt{[`G', `uk', `esh', `ĠDom', `mar', `aju']} \\
Mistral, DictaLM
  & \texttt{[`\spmark{}Gu', `k', `esh', `\spmark{}Dom', `mar', `aju']} \\
Granite
  & \texttt{[`G', `uk', `esh', `ĠDom', `mar', `aj', `u']} \\
SmolLM2
  & \texttt{[`G', `uk', `esh', `ĠD', `omm', `ar', `aj', `u']} \\
\bottomrule
\end{tabular}
\end{table}

Six of the ten models agree on a single 6-token segmentation.  The two
SentencePiece models (Mistral and DictaLM) produce the same alternative
6-token split, differing only in the boundary prefix character (\texttt{\spmark}
vs.\ \texttt{Ġ}), which reflects their different tokenizer type rather than
their training vocabulary.  Granite and SmolLM2 each produce unique
segmentations with 7 and 8 tokens respectively, indicating that the substring
\texttt{"Dommaraju"} was not frequent enough in their (smaller) training
corpora to be represented as a single merged token.

For the common phrase \textbf{\texttt{"The World Chess Championship"}}, the
majority (6 models: Llama, Qwen, OLMo, DeepSeek, Phi, Falcon) tokenize it into
4 tokens: \texttt{[`The', `ĠWorld', `ĠChess', `ĠChampionship']}.  The two
SentencePiece models split \texttt{"Chess"} as \texttt{[`\spmark{}Che', `ss']} (5
tokens), Granite fragmentizes \texttt{"Championship"} further into four pieces
(8 tokens total), and SmolLM2 splits \texttt{"Championship"} as
\texttt{[`ĠChampions', `hip']} (5 tokens).

\paragraph{Likely cause.}
The variation on proper names stems primarily from \emph{vocabulary size and
training corpus composition}.  Rare names that were seen infrequently receive
fewer merged subword units, causing them to be split into smaller fragments.
The boundary-marker divergence (\texttt{Ġ} vs.\ \texttt{\spmark}) is a consequence
of the underlying tokenizer algorithm (BPE vs.\ SentencePiece) and has no
effect on reconstruction fidelity.

% ----------------------------------------------------------
\subsection{Discussion}
% ----------------------------------------------------------

\paragraph{English efficiency.}
All models are similarly efficient for English, producing roughly 1.28--1.49
tokens per whitespace word.  The outlier is \textbf{Granite-3.3} (1.76), which
has the smallest base vocabulary (49{,}159 before special tokens) among
the BPE models and therefore covers fewer long English morphemes.

\paragraph{Hebrew efficiency.}
The results split the ten models into two clear groups:

\begin{itemize}[noitemsep]
  \item \textbf{Hebrew-aware} (HE/EN $\leq 3$): Qwen2.5, DeepSeek-V3, Phi-4-mini,
        DictaLM-2.0. These tokenizers have seen substantial Hebrew during
        training and have accumulated many whole-Hebrew-word tokens.
        Phi-4-mini achieves the best ratio (1.88$\times$) with 200{,}029
        vocabulary entries; DictaLM-2.0 benefits from being trained
        specifically on Hebrew/Israeli text.
  \item \textbf{Hebrew-inefficient} (HE/EN $> 4$): Llama, Mistral, OLMo,
        SmolLM2, Falcon, and Granite.  Hebrew words are fragmented into
        individual characters or very short byte sequences, producing 5--6
        tokens per word.  SmolLM2 is the worst at 6.39 tokens per Hebrew word
        ($4.80\times$ the English rate).
\end{itemize}

\paragraph{Implications for deployment.}
A high HE/EN ratio is costly: it inflates context length for Hebrew inputs,
reduces effective context window, and increases inference latency and cost.
For a Hebrew-facing application one should prefer Qwen2.5, Phi-4-mini,
DeepSeek-V3, or DictaLM-2.0.  For purely English tasks the difference is
negligible (all models within 0.5 tokens/word of each other).

% ============================================================
\newpage
\section{Part 3 - Constrained Decoding}
% ============================================================

% ----------------------------------------------------------
\subsection{Hebrew Token Identification}
% ----------------------------------------------------------

\paragraph{Strategy.}
For each token ID in a model's vocabulary we decode it in isolation using
\texttt{tokenizer.decode([id])} and apply the following rule:

\begin{quote}
A token is \emph{allowed} if and only if every character in its decoded string
that belongs to Unicode letter category (general category starting with
\texttt{L}) is also a Hebrew character, i.e.\ falls in the Unicode Hebrew block
U+0590--U+05FF.
\end{quote}

This rule keeps: Hebrew letter tokens, digit tokens, all punctuation and
whitespace tokens (their characters are never in category \texttt{L}),
byte-fallback tokens that decode as the replacement character U+FFFD (their
category is \texttt{So}), and special tokens such as EOS/BOS.
It rejects: any token containing a Latin, Cyrillic, Arabic, or other non-Hebrew
letter.
The EOS token is unconditionally added to the allowed set so that generation can
always terminate.

The token IDs are written to \texttt{hebrew\_allowed\_tokens\_qwen.json} and
\texttt{hebrew\_allowed\_tokens\_mistral.json}.

\paragraph{Results.}
Table~\ref{tab:hebrew-tokens} summarises the allowed-token counts for each model.

\begin{table}[H]
\centering
\begin{tabular}{lrrrr}
\toprule
Model & Full vocab & Allowed & Hebrew-bearing & Digits/punct/special \\
\midrule
Qwen2.5-7B-Instruct       & 151,665 & 13,234 (8.7\%) & 3,164 & 10,070 \\
Mistral-7B-Instruct-v0.3  &  32,768 &  1,912 (5.8\%) &    38 &  1,874 \\
\bottomrule
\end{tabular}
\caption{Hebrew-allowed token counts. ``Hebrew-bearing'' means the decoded token
contains at least one Hebrew character. The remainder are digits, punctuation,
whitespace, control, and special tokens that contain no letters.}
\label{tab:hebrew-tokens}
\end{table}

The large difference in Hebrew-bearing token counts (3,164 vs.\ 38) reflects
training data: Qwen2.5 was trained on multilingual corpora and has tokens for
common Hebrew letter sequences and words, while Mistral-7B-v0.3 was trained
predominantly on English and European text and retains only 38 Hebrew tokens ---
all single letters such as \textit{alef}, \textit{bet}, \textit{gimel}, etc.
Notably, Qwen's larger vocabulary also contributes to its much larger
\emph{non-Hebrew} allowed set (10,070 digits/punctuation/special tokens vs.\
1,874 for Mistral).

% ----------------------------------------------------------
\subsection{Implementation}
% ----------------------------------------------------------

Constrained decoding was implemented as a custom \texttt{LogitsProcessor} subclass,
\texttt{HebrewOnlyLogitsProcessor}, registered via HuggingFace's
\texttt{LogitsProcessorList} hook that is called immediately before each
sampling step.

\paragraph{Mechanism.}
On the first forward pass the processor builds a Boolean mask of length equal to
the actual logit tensor's vocabulary dimension (which can differ from
\texttt{tokenizer.vocab\_size} due to padding).
All positions corresponding to non-allowed token IDs are set to \texttt{False};
all positions in the allowed set are set to \texttt{True}.
At every decoding step the processor sets the logits of all masked-out tokens to
$-\infty$, making those tokens impossible to sample.
The mask is built lazily and cached, so the overhead is paid only once per
generation call.

\paragraph{Decoding configuration.}
Greedy decoding (\texttt{do\_sample=False}) was used throughout for
determinism.
Both constrained and unconstrained runs use identical prompts, formatted with
\texttt{tokenizer.apply\_chat\_template(messages, tokenize=False,}
\texttt{add\_generation\_prompt=True)}, which applies each model's native
conversational format (Qwen: \texttt{<|im\_start|>user\ldots<|im\_end|>};
Mistral: \texttt{[INST]\ldots[/INST]}).
This ensures a fair comparison: the only variable between the two conditions is
the presence or absence of the logits mask.

% ----------------------------------------------------------
\subsection{Results and Discussion}
% ----------------------------------------------------------

\paragraph{Output overview.}
The unconstrained outputs for all 10 queries are fluent, well-structured English
responses from both models. The full outputs appear in
\texttt{decoding\_outputs.jsonl}. Table~\ref{tab:constrained-outputs} shows the
constrained outputs, truncated to 120 characters.

\begin{longtable}{L{3.5cm} L{4cm} L{4cm}}
\toprule
Query & Qwen2.5-7B (constrained) & Mistral-7B (constrained) \\
\midrule
\endhead
\bottomrule
\endfoot
Why is the sky blue?
  & \texttt{!';{\textbackslash}n\#\#\# **[1. **[**[**[**[**[**[**[**[** ...}
  & \texttt{{\textbackslash}n{\textbackslash}n1. **(1) ** **** ** ** ** ** ...} \\[4pt]

Public transportation advantages/disadvantages
  & \texttt{\#\#\# **.);\textbackslash{}n- **(1)**\textbackslash{}n **(2)** ...}
  & \texttt{**1. **\textbackslash{}n\textbackslash{}n**-** **** **** **** ...} \\[4pt]

Email to professor for extension
  & \texttt{**[{\small(emoji)}]** [12/05/2023] **[{\small(emoji)}]** ...}
  & \texttt{[---]\textbackslash{}n\textbackslash{}n[**[**]**]\textbackslash{}n ...} \\[4pt]

How to make an omelette
  & \texttt{¡¡¡¡¡¡¡¡¡¡¡¡¡¡¡¡¡¡¡¡¡¡¡¡¡¡¡¡¡¡¡¡¡¡¡¡¡¡¡¡¡¡¡¡¡¡¡¡¡¡¡ ...}
  & \texttt{1. **1-2-3-4-5-6-7-8-9-10-11-12-13-14-15- ...} \\[4pt]

Supervised vs.\ unsupervised learning
  & \texttt{**1. ** **[** **\_** **\_** **\_** **\_** **\_** **\_** ...}
  & \texttt{**1. **\textbackslash{}n\textbackslash{}n**1.1. **\textbackslash{}n\textbackslash{}n**1.1.1. ** ...} \\[4pt]

Summarise Cinderella
  & \texttt{.`); //\{\} // (CJK char) // (CJK char) ...}
  & \texttt{"1) '[1]' ([2])'[3]' ([4])'[5]' ([6]) ...} \\[4pt]

Reduce smartphone distraction
  & \texttt{1. **{\small(emoji)}** "{\small(emoji)}" "{\small(emoji)}" ...}
  & \texttt{1. **"**\_"**\_**\_**\_**\_**\_**\_**\_**\_ ...} \\[4pt]

What happens when water boils
  & \texttt{ !"); //\textbackslash{}n // 3.1.2.2.2.2.2.2.2.2.2. ...}
  & \texttt{1. **100°C (212°F)**:\textbackslash{}n\textbackslash{}n - **100°C** (760.00714 ...} \\[4pt]

Polite refusal to party invitation
  & \texttt{¡¡¡¡¡¡¡¡¡¡¡¡¡¡¡¡¡¡¡¡¡¡¡¡¡¡¡¡¡¡¡¡¡¡¡¡¡¡¡¡¡¡¡¡¡¡¡ ...}
  & \texttt{"**[**\_[**]**],\textbackslash{}n\textbackslash{}n**[**\_[**]**],** ...} \\[4pt]

``Practice makes progress'' advice
  & \texttt{" ,'"\textbackslash{}n\textbackslash{}n" ,'"\textbackslash{}n\textbackslash{}n" ,'"\textbackslash{}n\textbackslash{}n" ,'" ...}
  & \texttt{"**'**\_'**\_'**\_'**\_'**\_'**\_'**\_'**\_ ...} \\
\caption{Constrained decoding outputs (truncated). Full outputs in \texttt{decoding\_outputs.jsonl}.}
\label{tab:constrained-outputs}
\end{longtable}

\paragraph{Is the Hebrew text meaningful?}
No Hebrew characters appear in any constrained output from either model.
Despite Hebrew-bearing tokens being present in both allowed sets (3,164 for Qwen;
38 for Mistral), greedy decoding never selects them.

\paragraph{Does the output relate to the question?}
No. The constrained outputs contain no semantic content; they are purely
structural or punctuation-based symbol sequences with no informational value.

\paragraph{Do the two models differ?}
The outputs are gibberish for both models, but they follow different degenerate
patterns reflecting each model's internal ``schema memory'':

\begin{itemize}[noitemsep]
\item \textbf{Qwen} tends to produce Markdown-like header loops
  (\texttt{\#\#\# **[1. **[**[**[** \ldots}) or repeating single symbols
  (\texttt{¡¡¡¡¡ \ldots}, \texttt{" ,'" \ldots}).
\item \textbf{Mistral} tends to produce numbered list or hierarchical heading
  loops (\texttt{**1. **}, \texttt{**1.1. **}, \texttt{**1.1.1. **}, \ldots),
  template bracket structures (\texttt{[**[**]**] \ldots}), or
  incrementing number sequences (\texttt{1-2-3-4-5- \ldots}).
\end{itemize}

Both patterns suggest the model enters a degenerate attractor state once normal
text tokens are unavailable.

\paragraph{Are some languages favoured over others?}
No natural language is produced. The outputs are language-free sequences of
punctuation, digits, and structural tokens.

\paragraph{Root cause.}
The core problem is a \emph{mismatch between the token constraint and the
model's probability distribution}.
When the model is conditioned on an English prompt, its distribution assigns
probabilities close to zero to Hebrew tokens and high probabilities to
English-letter tokens.
Masking out all non-Hebrew-allowed tokens removes the high-probability English
tokens, but leaves intact a large set of punctuation and structural tokens that
the model assigns moderate probability to.
Greedy decoding therefore collapses to whichever punctuation token has the
highest logit at each step, quickly entering a repetitive loop.

In principle, adding a Hebrew-language instruction to the prompt (e.g.,
``Please respond in Hebrew'') would shift the distribution toward Hebrew tokens
and likely produce meaningful Hebrew output under the same constraint.
However, this was deliberately not done: modifying the prompt for the
constrained condition but not the unconstrained condition would make the two
conditions incomparable, since any quality difference could be attributed to the
prompt change rather than to the token constraint alone.
Fine-tuning the model to shift its \emph{prior} distribution toward Hebrew ---
the subject of Part~4 --- is the correct approach.

% ============================================================
\newpage
\section{Part 4 - Fine Tuning}
% ============================================================

% ----------------------------------------------------------
\subsection{Data}
% ----------------------------------------------------------

For fine-tuning, I created a supervised instruction dataset of
\textbf{1,000 English prompt $\to$ Hebrew answer pairs}.
The prompts were sampled from the
\texttt{tatsu-lab/alpaca} instruction corpus, and the Hebrew responses were
generated automatically using the Gemini API in
\texttt{code/part4/create\_data.py}.
The script uses \texttt{gemini-3-flash-preview} with the system instruction
``Always answer in Hebrew'', then stores each pair in JSONL format.

To prevent data leakage, the 20 assignment evaluation prompts are explicitly
listed in the script (\texttt{EVAL\_INPUTS}) and filtered out before
writing training examples. Therefore, none of the evaluation inputs appears in
the fine-tuning set.

The final training set covers diverse topics (science, history, writing,
everyday reasoning, and basic programming), which helps evaluate whether the
model learned a general language behavior (answering in Hebrew) rather than a
single narrow template.

% ----------------------------------------------------------
\subsection{Fine-Tuning Procedure}
% ----------------------------------------------------------

I fine-tuned \texttt{Qwen/Qwen2.5-1.5B-Instruct} using LoRA (PEFT) with
\texttt{trl} SFTTrainer in Google Colab.
LoRA was selected to reduce compute and memory cost while keeping base model
weights frozen.

\begin{table}[H]
\centering
\caption{LoRA Configuration}
\begin{tabular}{L{4cm}L{8.5cm}}
\toprule
Parameter & Value \\
\midrule
Rank $r$ & 8 \\
$\alpha$ & 16 \\
Dropout & 0.05 \\
Target modules & \texttt{q\_proj}, \texttt{k\_proj}, \texttt{v\_proj}, \texttt{o\_proj} \\
\bottomrule
\end{tabular}
\end{table}

\begin{table}[H]
\centering
\caption{Training Hyper-Parameters}
\begin{tabular}{L{4cm}L{8.5cm}}
\toprule
Parameter & Value \\
\midrule
Epochs & 3 \\
Batch size & 2 \\
Gradient accumulation & 4 (effective batch size = 8) \\
Learning rate & $2 \times 10^{-4}$ \\
Max sequence length & 512 \\
\bottomrule
\end{tabular}
\end{table}

I used a completion-only data collator, so loss is computed only on assistant
tokens (Hebrew response) while user-prompt tokens are masked with label $-100$.
This directly optimizes the target behavior: produce a Hebrew response for an
English instruction.

The same Qwen chat template format was used consistently in both training and
inference:
\texttt{<|im\_start|>user ... <|im\_end|><|im\_start|>assistant ...}.

% ----------------------------------------------------------
\subsection{Evaluation}
% ----------------------------------------------------------

I evaluated the original and fine-tuned model on all 20 required prompts
(stored in \texttt{eval\_outputs.jsonl}).

\begin{table}[H]
\centering
\caption{Evaluation Summary (20 prompts)}
\begin{tabular}{L{5cm}C{4cm}C{4cm}}
\toprule
Metric & Base model & Fine-tuned model \\
\midrule
Responses in Hebrew & 0/20 (0\%) & 20/20 (100\%) \\
Good / Partial / Poor relevance & 20 / 0 / 0 & 3 / 11 / 6 \\
\bottomrule
\end{tabular}
\end{table}

\textbf{Training-loss evidence of learning.}
The training loss started at \textbf{2.825} and decreased to
\textbf{2.355}.
In this setup, loss is the average token-level negative log-likelihood on the
assistant response tokens:
\[
\mathcal{L} = -\frac{1}{N}\sum_{t=1}^{N}\log p_{\theta}(y_t\mid x, y_{<t}).
\]
A lower value means the model assigns higher probability to the correct next
Hebrew token, so the downward trend indicates real learning of the target
generation behavior.

\textbf{Two good examples and one bad example.}
\begin{itemize}[leftmargin=*]
  \item \textbf{Good: ``What is the capital of France?''}
  The fine-tuned model answered correctly in Hebrew and identified Paris.
  \item \textbf{Good: ``Write a short poem about rain.''}
  The fine-tuned model produced a coherent Hebrew poem structure, showing it
  learned the requested output language and style format.
  \item \textbf{Bad: ``Name two famous painters from the Renaissance.''}
  The fine-tuned model hallucinated and listed Shakespeare as a painter,
  which is factually incorrect.
\end{itemize}

\textbf{Per-prompt results (all 20 inputs).}
\begin{longtable}{L{0.8cm}L{6.2cm}C{2.2cm}C{2.2cm}L{3.0cm}}
\caption{Base vs Fine-Tuned Outputs on the 20 Evaluation Prompts}\label{tab:part4_all_outputs}\\
\toprule
ID & Prompt (short) & Base lang & FT lang & FT quality note \\
\midrule
\endfirsthead
\toprule
ID & Prompt (short) & Base lang & FT lang & FT quality note \\
\midrule
\endhead
1 & Why sky looks blue & English & Hebrew & Poor (topic drift) \\
2 & 2 pros/2 cons public transport & English & Hebrew & Partial \\
3 & Email for extension & English & Hebrew & Partial \\
4 & How to make omelette & English & Hebrew & Poor \\
5 & Supervised vs unsupervised & English & Hebrew & Partial \\
6 & Cinderella in 3 sentences & English & Hebrew & Poor \\
7 & Reduce smartphone distraction & English & Hebrew & Partial \\
8 & What happens when water boils & English & Hebrew & Partial \\
9 & Polite refusal to party & English & Hebrew & Partial \\
10 & Practice makes progress advice & English & Hebrew & Poor \\
11 & Capital of France & English & Hebrew & Good \\
12 & Number of continents & English & Hebrew & Good \\
13 & Gravity in simple terms & English & Hebrew & Partial \\
14 & Short poem about rain & English & Hebrew & Good \\
15 & 3 exercise benefits & English & Hebrew & Partial \\
16 & How to make tea & English & Hebrew & Partial \\
17 & Noun vs verb difference & English & Hebrew & Partial \\
18 & Renaissance painters & English & Hebrew & Poor (hallucination) \\
19 & Why people dream & English & Hebrew & Partial \\
20 & Advice for a new job & English & Hebrew & Partial \\
\bottomrule
\end{longtable}

Overall, fine-tuning was \textbf{successful for language control}
(100\% Hebrew outputs) but only \textbf{partially successful for factual and
instruction accuracy}.
The model appears to have learned a strong Hebrew-response prior, while content
quality remains mixed on more complex prompts.

\end{document}
